\begin{table}[H]
\centering
\caption{Erro de Calibração do Value at Risk — Validação}
\label{tab:var}
\begin{tabular}{lcccccc}
\hline
 & \multicolumn{2}{c}{\textbf{q = 1\%}} & \multicolumn{2}{c}{\textbf{q = 5\%}} & \multicolumn{2}{c}{\textbf{q = 10\%}} \\
\cmidrule(lr){2-3} \cmidrule(lr){4-5} \cmidrule(lr){6-7}
\textbf{Modelo} & \textbf{Erro} & \textbf{Emp.} & \textbf{Erro} & \textbf{Emp.} & \textbf{Erro} & \textbf{Emp.} \\
\hline
NeuralFactors & 0,4156 & 0,4256 & 0,3871 & 0,4371 & 0,3493 & 0,4493 \\
PPCA          & 0,0094 & 0,0194 & 0,0048 & 0,0452 & 0,0216 & 0,0784 \\
\hline
\end{tabular}
\caption*{\footnotesize \textbf{Nota:} Os erros de calibração do NeuralFactors apresentados nesta tabela são manifestamente inconsistentes com o desempenho esperado do modelo — erros da magnitude observada (0,35–0,42) indicam falha no pipeline de dados e não refletem a qualidade probabilística real do modelo. A título de referência, um modelo bem calibrado deveria apresentar erros próximos de zero para todos os quantis. Esses resultados serão revisados na versão final do trabalho.}
\caption*{Fonte: Elaboração própria.}
\end{table}